\section{Scope of This Version and Distribution}

\subsection{Stages}

\texttt{fsb} is delivered in three stages that grow the interface without changing
the security core: \emph{crawl} (browse), \emph{walk} (preview and search) and
\emph{run} (more).  This version implements the crawl and walk stages, as specified
in Sections~\ref{sec:invocation} to~\ref{sec:http}.  It has not been run against real
cloud-placeholder files or in Firefox; those checks are still to be done.

\subsection{Not in this version}

The following are wanted but are not part of this version.  They are recorded here so
that they are not designed away.

\begin{req}{FUT-1}
Full-text and fuzzy search by delegating to the system's Spotlight index, with every
result checked against the deny and ignore rules before it is shown.
\end{req}

\begin{req}{FUT-2}
Live updates when files change.
\end{req}

\begin{req}{FUT-3}
Several panes, tabs, and saved views.
\end{req}

\begin{req}{FUT-4}
Version-control status overlays.
\end{req}

\begin{req}{FUT-5}
``Reveal in Finder'' and ``Open in editor'', restricted to a fixed list of actions,
never building a command from a path.  Excluded from this version because they would
make the server start local processes at the user's request; ``Copy path'' covers the
need meanwhile.  The one process this version starts is Quick Look's \code{qlmanage},
a fixed program run on a private copy (\rid{SEC-15}).
\end{req}

\begin{req}{FUT-6}
Derived data (an index, thumbnails, caches), if ever added, shall honor the deny
rules when it is built and when it is queried.
\end{req}

\subsection{Distribution}

\begin{req}{DST-1}
The program shall be published in a public repository under the MIT license.
\end{req}

\begin{req}{DST-2}
Releases are source releases: a version tag in the public repository, installed with
\code{go install github.com/rprimmer/fsb/cmd/fsb@latest} or built with \code{make
install} from a clone.  No prebuilt binaries are provided, so none needs signing or
notarization; a program built locally is not quarantined by macOS.
\end{req}

\begin{req}{DST-5}
The source shall provide \code{make install}, which installs the program and its manual
page; the installation locations (\code{PREFIX}, \code{BINDIR}, \code{MANDIR}) shall be
variables that the user can override, and \code{DESTDIR} shall stage the installation for
packaging.  The manual page goes in the \code{man1} folder of \code{MANDIR}.
\end{req}

\begin{req}{DST-3}
No Homebrew formula is provided.  (Planned earlier; withdrawn, because maintaining
one is ongoing work for a project that is finished.)
\end{req}

\begin{req}{DST-4}
The documentation shall explain the security model, the deny rules and their costs,
and that macOS may ask for permission to read Documents, Desktop, Downloads and
external volumes the first time.
\end{req}

Version 1.0 is released as source (\rid{DST-2}), and the project is provided as-is: it
is finished and not actively maintained.
